\documentclass[11pt]{article}

\usepackage[margin=1in]{geometry}
\usepackage{amsmath, amssymb, amsthm}
\usepackage{algorithm}
\usepackage{algpseudocode}
\usepackage{graphicx}
\usepackage{hyperref}
\usepackage[utf8]{inputenc}
\usepackage[T1]{fontenc}
\usepackage{booktabs}
\usepackage{caption}
\usepackage{float}
\usepackage{enumitem}

\hypersetup{
    colorlinks=true,
    linkcolor=blue,
    citecolor=blue,
    urlcolor=blue
}

\title{
    \textbf{\LARGE WiFi Access Point Placement}\\
    \textbf{\LARGE Using Gradient Descent with Backtracking Line Search}\\[0.2cm]
    \large Numerical Optimization - Project Report
}

\author{
    \large Anđela Maksimović 20009
}

\date{\today}

\begin{document}

\maketitle

\begin{abstract}
This report presents a numerical optimization approach for placing a
fixed number of wireless access points (APs) inside a building in order
to maximize signal coverage. The problem is formulated as a constrained
nonlinear optimization problem, where the decision variables are the
coordinates of the APs. A projected gradient ascent method is used to
solve the problem, with the gradient estimated using central finite
differences and the step size selected using a backtracking line
search. The building is represented as a two-dimensional grid, allowing
walls and obstacles to be included through signal attenuation. An
interactive implementation is also provided to visualize the floor
plan, AP positions, coverage, and optimization process. The main goal
of the project is to demonstrate the application and behavior of a
gradient-based numerical optimization method on a practical,
non-convex problem.
\end{abstract}

\tableofcontents

% ==================================================================
\section{Introduction and Background}
% ==================================================================

The placement of wireless access points is an important part of
designing a WiFi network. The locations of the access points affect the
area that can be covered and the quality of the received signal. When a
building contains multiple rooms, corridors, and obstacles, finding
suitable AP locations manually becomes more difficult.

The problem can be viewed as an optimization problem in which the
positions of a fixed number of APs are changed in order to maximize a
coverage measure. Similar wireless network optimization problems have
been studied using different objectives and optimization methods.
Coverage and signal propagation are among the most common objectives,
while other approaches consider throughput and load balancing
\cite{mcgibney2007user}, channel assignment
\cite{lee2002optimization}, or localization accuracy
\cite{battiti2003optimal}.

From an optimization perspective, AP placement can be approached in
several ways. Gradient-based methods use information about how the
objective changes when the decision variables are modified; for
example, differentiable formulations of cellular network throughput
have been optimized directly through gradient descent
\cite{eller2024differentiable}. Classical methods such as gradient
descent, gradient ascent, quasi-Newton, and conjugate-gradient methods
belong to this group. Other approaches are more suitable when the
problem is discrete or when gradient information is not available.
These include direct-search methods and metaheuristics such as
genetic algorithms, simulated annealing, and tabu search \cite{molina1999automatic}.

For this project, a gradient-based approach was selected because it
provides a direct way to study several concepts from numerical
optimization. In particular, the project considers numerical gradient
estimation, constrained optimization through projection, and
backtracking line search for selecting the step size. Since the
resulting problem is non-convex, the algorithm is expected to find a
local optimum rather than necessarily the global optimum.

\subsection{Related implementations}

Several existing implementations were considered when designing the
problem and its visualization.

\begin{itemize}
    \item \textbf{WSN-Optimization-using-GA} uses a genetic algorithm for wireless sensor network placement and
    provides a useful comparison with a gradient-based approach \cite{wsn_ga}.

    \item \textbf{plexus-network-planner}
    was used as inspiration for the interactive floor-plan
    visualization, particularly the idea of displaying AP positions
    and coverage directly on a building layout \cite{plexus_planner}.

    \item \textbf{Wireless-Access-Point-Optimization}
    was considered as a reference for AP placement and coverage
    calculation \cite{wap_optimization}.
\end{itemize}

These projects were used as references for the general problem and
visualization. The optimization procedure used in this project was
implemented separately.

% ==================================================================
\section{Problem Statement}
% ==================================================================

The objective of the project is to determine the placement of a fixed
number of wireless access points inside a building in order to maximize
signal coverage.

The decision variables are the coordinates of the APs. For $n$ access
points, the optimization variable is

\[
u=(x_1,y_1,\ldots,x_n,y_n),
\]

where $(x_i,y_i)$ represents the position of access point $i$.

The problem is constrained because an AP cannot be placed outside the
building or inside a wall or obstacle. The resulting optimization
landscape is non-convex and may contain multiple local optima. This
means that different initial AP configurations can lead to different
final solutions.

% ==================================================================
\section{Optimization Model}
% ==================================================================

\subsection{Floor-plan representation}

The building is represented using a two-dimensional grid. Each grid
cell describes a part of the environment and is classified as
\texttt{Outside}, \texttt{Free}, \texttt{Wall}, or \texttt{Obstacle}.
Only \texttt{Free} cells are valid positions for an AP.

This representation allows the building layout to be used directly by
the optimization procedure. Different floor plans can be generated
with different spatial structures, including open rooms, corridors,
L-shaped layouts, and multi-room office-complex layouts (kitchen,
meeting room, and open workspace connected through partition
doorways). This makes it possible to test the optimization method
under different constraints and non-convex geometries.

Each cell also contains an attenuation coefficient
$\alpha(g_x,g_y)\geq0$, which represents the effect of walls or other
obstacles on the signal.



\subsection{Coverage objective}

The signal contributions from multiple APs are combined to obtain the
coverage of a grid point. Assuming independent coverage contributions,
the probability that a point $q$ is covered by at least one AP is

\[
C(u,q)=1-\prod_{i=1}^{n}(1-s_i(u,q)).
\]

The total coverage over all free cells $\mathcal{Q}$ is then

\[
g(u)=\sum_{q\in\mathcal{Q}}C(u,q).
\]

To discourage APs from being placed too close to each other, a
repulsion penalty is added. The minimum desired distance between two
APs is defined as

\[
d_{\min}=0.8R(n).
\]

The final objective function is

\[
f(u)=g(u)
-\sum_{i<j}
2\big(d_{\min}-\lVert u_i-u_j\rVert_2\big)^2
\mathbb{1}\!\left[
0<\lVert u_i-u_j\rVert_2<d_{\min}
\right].
\]

The optimization problem is therefore

\[
\max_{u\in\mathcal{F}} f(u),
\]

where $\mathcal{F}$ represents all valid AP placements inside the
building.

The objective is non-convex because of the nonlinear signal model,
coverage combination, and repulsion term. In addition, the feasible
region created by the floor plan is generally non-convex and can
contain disconnected areas. As a result, the optimization method is
not expected to provide a guarantee of a global optimum.

% ==================================================================
\section{Numerical Method}
% ==================================================================
\subsection{Signal model}

Since AP and evaluation points are continuous but the floor plan is
discretized on a grid, the set of cells crossed by a line-of-sight ray
is found using Bresenham's line-rasterization algorithm
\cite{bresenham1965algorithm}, the standard technique for this problem
in computer graphics. The total attenuation along
this line is calculated as

\[
A(p,q)=
\sum_{(g_x,g_y)\in\text{Bresenham}(p,q)}
\alpha(g_x,g_y).
\]

This provides a simple way of taking walls and obstacles into account
when calculating the signal received at a particular grid cell.

Each AP transmits at a fixed power, and the total power budget
$P_{\text{tot}}$ available to the network is shared equally between
the $n$ APs, so each AP transmits at $P_{\text{tot}}/n$. The effective
range of a single AP is defined relative to a reference configuration:
$P_0$ is the reference transmit power of one AP, and $R_0$ is the
range that AP would achieve at that reference power. Since received
power decreases with the square of distance, the range an AP can reach
scales with the square root of its transmit power relative to this
reference, so the effective range when $n$ APs share the power budget
is

\[
R(n)=R_0\sqrt{\dfrac{P_{\text{tot}}/n}{P_0}}.
\]

Increasing the number of APs therefore reduces the effective range of
each individual AP, reflecting the fact that the same total power is
now split among more transmitters.

For an AP at position $u_i$ and a grid point $q$, the distance is

\[
d_i=\lVert q-u_i\rVert_2.
\]

The signal contribution of AP $i$ is modeled as

\[
s_i(u,q)=
\frac{1}{1+\exp\!\big(k(d_i-0.75R(n))\big)}
e^{-A(u_i,q)},
\qquad
k=\frac{6}{R(n)}.
\]

The first part of the expression models the decrease in signal with
distance, while the second part reduces the signal according to the
attenuation caused by the floor-plan obstacles.

\subsection{Projection onto the feasible region}

A normal gradient step can move an AP into a wall or outside the
building. To handle this constraint, the updated position is projected
onto the feasible region after every gradient step.

The projection is implemented using the floor-plan grid. Coordinates
are first restricted to the floor-plan boundaries. If the resulting
point is already inside a \texttt{Free} cell, it is kept unchanged.
Otherwise, nearby cells are searched in expanding Chebyshev rings until
a \texttt{Free} cell is found. The AP is then moved to the center of
that cell.

This is an approximate projection rather than an exact Euclidean
projection, but it provides a simple way to keep every accepted AP
position inside the valid part of the building.

\subsection{Backtracking line search}

The gradient provides the direction of the update, but an appropriate
step size still has to be selected. A backtracking line search is used
for this purpose.

Starting with an initial step size $\alpha_0$, the step is repeatedly
reduced according to

\[
\alpha\leftarrow0.7\alpha
\]

until the projected candidate satisfies the sufficient-increase
condition

\[
u^+(\alpha)=
\Pi_{\mathcal{F}}
\big(u+\alpha\nabla f(u)\big),
\]

\[
f(u^+(\alpha))
\geq
f(u)+c_1\alpha\lVert\nabla f(u)\rVert_2^2.
\]

The projection is performed before evaluating the condition. Therefore,
a candidate is only accepted after it has been moved back into the
feasible region.

\subsection{Stopping criteria}

The optimization stops when one of the following conditions is
satisfied:

\begin{itemize}
    \item the gradient norm becomes smaller than $10^{-4}$;
    \item the improvement in the objective becomes smaller than
    $10^{-6}$;
    \item no suitable step size can be found by the line search; or
    \item the maximum number of 200 iterations is reached.
\end{itemize}

The objective value and step size are recorded during the optimization
so that the convergence behavior can be inspected afterwards.

\begin{algorithm}[H]
\caption{Projected Gradient Ascent with Backtracking Line Search}
\label{alg:gradient}
\begin{algorithmic}[1]
    \State $u\gets\Pi_{\mathcal F}(u_0)$
    \State $v\gets f(u)$

    \For{$k=0,\dots,K_{\max}-1$}
        \State $\nabla\gets
        \text{CentralDifferenceGradient}(f,u,\varepsilon)$

        \If{$\lVert\nabla\rVert_2<\tau_g$}
            \State \textbf{stop} (GradientNorm)
        \EndIf

        \State $\alpha\gets\alpha_0$
        \State found $\gets$ false

        \For{$b=1,\dots,B_{\max}$}
            \State $u^+\gets
            \Pi_{\mathcal F}(u+\alpha\nabla)$
            \State $v^+\gets f(u^+)$

            \If{$v^+\geq
            v+c_1\alpha\lVert\nabla\rVert_2^2$}
                \State found $\gets$ true
                \State \textbf{break}
            \EndIf

            \State $\alpha\gets0.7\alpha$
        \EndFor

        \If{not found}
            \State \textbf{stop} (LineSearchFailed)
        \EndIf

        \State $u\gets u^+$
        \State $v\gets v^+$

        \If{$v^+-v<\tau_f$}
            \State \textbf{stop} (ObjectiveStall)
        \EndIf
    \EndFor

    \State \Return $u$, $v$, iteration count, stopping reason,
    histories
\end{algorithmic}
\end{algorithm}

\subsection{Computational cost}

Let $|\mathcal{Q}|$ denote the number of free cells and $L$ the average
length of a rasterized line between an AP and a grid cell. One
evaluation of the coverage objective requires approximately

\[
O(n|\mathcal{Q}|L)
\]

operations for the signal calculations and $O(n^2)$ operations for the
repulsion term.

There are $2n$ decision variables, and central finite differences
require two objective evaluations for each variable. Therefore, one
gradient calculation requires approximately $4n$ objective
evaluations. The computational cost increases with both the number of
APs and the resolution of the floor-plan grid.

% ==================================================================
\section{Implementation}
% ==================================================================

The project separates the numerical optimization core from the
graphical interface used to visualize it. The floor plan, signal and
coverage calculation, gradient estimation, projection, and
optimization loop are implemented independently of any GUI code and
operate purely on dense matrices (natID) and coordinate data. The interactive
visualization layer, built with the \texttt{natid} GUI framework, only
wraps this core: it displays the floor plan, current AP positions, and
coverage information, and drives the optimizer through the same
public interface used by any other caller. This separation means the
optimization procedure can be exercised, tested, or reused
independently of the GUI.

The floor plan is generated as a two-dimensional grid, where each cell
contains information about whether it can be occupied and, when
appropriate, its attenuation value. The generator can create several
structured layouts, such as open rooms, corridors, L-shaped
environments, and multi-room office-complex layouts. These different
layouts provide benchmark instances with different feasible regions
and obstacle configurations.

The generated floor plan is then passed directly to the optimization
pipeline. This allows the same optimization method to be tested under
different spatial constraints without changing the numerical
procedure.

For the grid representation, the occupancy and attenuation data are
stored using dense matrix structures. This also makes it possible to
save and load the attenuation data in MatrixMarket format when needed
during development or analysis.

The interactive part of the application displays the floor plan,
current AP positions, and coverage information. The optimization
history can also be used to observe how the objective and step size
change during the iterations. The GUI is therefore mainly used to make
the numerical method easier to inspect and understand.

% ==================================================================
\section{Optimization Behavior}
% ==================================================================

Because the optimization problem is non-convex, the starting positions
of the APs can influence the final result. Starting the algorithm from
different feasible configurations can lead to different local optima.

This behavior can be observed using the interactive implementation.
The AP positions can be randomized and the optimization can then be
run again. Comparing different runs makes it possible to see how the
initial configuration affects the final AP positions and the obtained
coverage.

The visualization is particularly useful for understanding the
optimization process because the changes are not limited to the final
objective value. The APs can be observed as they move through the
feasible region, while the objective and step-size histories provide
additional numerical information about the iterations.

This also illustrates an important limitation of the chosen method.
Projected gradient ascent uses local gradient information, so it does
not guarantee that the global maximum will be found. For a non-convex
problem such as this one, different initial configurations can result
in different solutions.

% ==================================================================
\section{Conclusion}
% ==================================================================

This project formulates WiFi access point placement as a constrained,
nonlinear optimization problem in which the coordinates of a fixed
number of APs are optimized to maximize signal coverage. The building
is represented as a two-dimensional grid, allowing the optimization to
take walls, obstacles, and different floor-plan geometries into
account.

The problem is solved using projected gradient ascent with a numerical
gradient and backtracking line search. Since the feasible region is
defined by the floor-plan grid, a grid-based projection is used to keep
the APs inside valid locations. The backtracking line search adapts the
step size during the optimization, while the stopping criteria provide
several ways of detecting convergence or termination.

The implementation also provides an interactive way to observe the
behavior of the algorithm. In particular, changing the initial AP
positions demonstrates the effect of non-convexity and the possibility
of converging to different local optima.

Overall, the project demonstrates the use of a gradient-based numerical
optimization method on a practical constrained problem. The main focus
is not on finding a universally optimal WiFi placement strategy, but on
showing how gradient estimation, projection, line search, and
convergence criteria can be combined to solve and analyze a
non-convex optimization problem.

The application is cross-platform and runs natively on macOS, Linux,
and Windows. Pre-built installers for all three platforms are
published as GitHub releases\footnote{\url{https://github.com/maksim0vic/ProjNO_WAPPlacement_Maksimovic/releases}},
allowing the tool to be tried directly without building it from
source.

\bibliographystyle{unsrt}  
\bibliography{references}  

\end{document}